In this work, we have demonstrated the theoretical validity and computational potential of a Spatial Photonic Ising Machine. By mapping the NP-hard ground-state search effectively onto the energy minimization of an optical interference pattern, we circumvent the von Neumann bottleneck that plagues conventional electronic solvers. Our simulation results confirm that the system correctly implements the intended Hamiltonian dynamics, exhibiting ferromagnetic ordering and spontaneous symmetry breaking consistent with mean-field theory.

The scalability of the SPIM architecture, limited primarily by the pixel density of commercial SLMs rather than thermal management, positions it as a superior candidate for solving large-scale combinatorial optimization problems. We have outlined concrete pathways for deployment in logistics, finance, and drug discovery, sectors where the ability to optimize complex systems translates directly to economic value.

Future work will focus on the experimental validation of these simulations using next-generation phase-only SLMs and the development of hybrid algorithms that couple the global search capability of the photonic machine with local gradient-based refinement on classical processors. As the physical limits of Moore's Law induce a plateau in digital scaling, the optical domain offers a vast, untapped coherent state space for the next era of high-performance computing.

\section*{Code Availability}
The simulation framework developed for this study, including an optimized Python library and an interactive web-based demonstration, is open-source and available at \url{https://github.com/alanspace/Photonic_Computing}. We encourage the community to adapt and extend this toolset for further research into optical computing architectures.
